% ========================================
% CHAPITRE 2 - PROBLÉMATIQUE ET CONTEXTE DU PROJET
% ========================================

\chapter{Probl\'ematique et Contexte du Projet}
\label{chap:problematique_contexte}

\section*{Introduction}

Ce chapitre pr\'esente de mani\`ere d\'etaill\'ee la probl\'ematique rencontr\'ee dans la gestion commerciale des offres t\'el\'ecom au sein de l'entreprise cliente, ainsi que le contexte sp\'ecifique qui a motiv\'e la conception et le d\'eveloppement de la plateforme Bifr\"ost. \`A travers une analyse approfondie des processus existants, nous identifierons les d\'efis majeurs, les points de douleur et les opportunit\'es d'am\'elioration qui ont justifi\'e la n\'ecessit\'e d'une solution sur mesure, moderne et int\'egr\'ee.

Cette analyse permet de comprendre les enjeux op\'erationnels, techniques et strat\'egiques qui sous-tendent le projet, tout en d\'efinissant clairement les objectifs \`a atteindre et les b\'en\'efices attendus de la mise en place d'une plateforme d\'edi\'ee \`a la gestion des devis et des abonnements B2B.

% ========================================
\section{Analyse des Processus Existants}
% ========================================

\subsection{Vue d'Ensemble des Processus Actuels}

Avant le d\'eveloppement de Bifr\"ost, la gestion des devis commerciaux au sein de l'entreprise cliente reposait sur des processus partiellement automatis\'es, intimement li\'es \`a un ERP g\'en\'eraliste inadapt\'e aux sp\'ecificit\'es m\'etier du secteur des t\'el\'ecommunications. Ces processus couvraient diff\'erentes \'etapes du cycle commercial~: la cr\'eation d'une opportunit\'e, la r\'edaction du devis, la n\'egociation commerciale, la validation et enfin la g\'en\'eration du bon de commande. Chacune de ces \'etapes mobilisait plusieurs acteurs et g\'en\'erait des interactions manuelles co\^uteuses en temps et en fiabilit\'e.

\subsubsection{Acteurs Impliqu\'es dans les Processus}

La cha\^ine de traitement des devis impliquait plusieurs profils d'utilisateurs aux r\^oles et responsabilit\'es distincts~:

\begin{itemize}
    \item \textbf{Agents commerciaux}~: responsables de la cr\'eation et de l'\'edition des devis \`a partir du catalogue d'offres.
    \item \textbf{Superviseurs ADV}~: charg\'es de la validation des devis avant transmission au client.
    \item \textbf{\'Equipes de production}~: destinataires des lignes de devis apr\`es validation, en charge de la mise en \oe{}uvre des services vendus.
    \item \textbf{Syst\`emes tiers (CRM, ERP)}~: syst\`emes sources et cibles des donn\'ees de devis, n\'ecessitant des \'echanges structur\'es et fiables.
\end{itemize}

\subsubsection{Flux du Processus Commercial Existant}

Le processus commercial suivait un enchaînement s\'equentiel rigide~:

\begin{enumerate}
    \item \textbf{Cr\'eation de l'opportunit\'e}~: le CRM g\'en\`ere une opportunit\'e commerciale et la transmet au syst\`eme de devis.
    \item \textbf{S\'election des articles}~: le commercial parcourt manuellement un catalogue de pr\`es de 45\,000 r\'ef\'erences pour s\'electionner les articles un par un.
    \item \textbf{\'Edition du devis}~: les lignes sont ajout\'ees, configur\'ees et valid\'ees manuellement, sans assistance de configurateur.
    \item \textbf{Validation}~: le superviseur ADV examine et valide le devis, sans interface d\'edi\'ee ni workflow automatis\'e.
    \item \textbf{Transmission}~: le devis valid\'e est transmis au client et aux \'equipes de production par des canaux non unifi\'es.
    \item \textbf{Archivage}~: le devis est archiv\'e manuellement lorsque son \'etat passe \`a inactif dans le syst\`eme source.
\end{enumerate}

\subsection{Processus de Gestion des Devis}

\subsubsection{\'Etape de Cr\'eation et de Configuration}

La phase de cr\'eation d'un devis constituait l'\'etape la plus chronophage et la plus source d'erreurs du processus existant. L'agent commercial devait~:

\begin{enumerate}
    \item Rechercher les articles adapt\'es parmi un catalogue de plusieurs dizaines de milliers de r\'ef\'erences, sans filtrage intelligents ni assistant de configuration~;
    \item Saisir manuellement chaque ligne avec ses param\`etres (quantit\'e, prix, remise, dur\'ee d'engagement)~;
    \item G\'erer manuellement la hi\'erarchie lignes\,/\,sous-lignes pour les offres complexes~;
    \item V\'erifier la coh\'erence des donn\'ees sans outil de validation automatis\'e.
\end{enumerate}

\subsubsection{\'Etape de Validation et de Transmission}

La validation des devis s'op\'erait de mani\`ere informelle, sans workflow d\'efini~:

\begin{enumerate}
    \item Le commercial soumettait le devis au superviseur ADV par email ou via un outil g\'en\'eraliste~;
    \item Le superviseur effectuait une relecture manuelle sans vue consolid\'ee des modifications~;
    \item Les retours \'etaient communiqu\'es de mani\`ere non structur\'ee, rallongeant les cycles de validation~;
    \item La transmission aux \'equipes de production se faisait sans correspondance automatique entre les lignes commerciales et les lignes de production.
\end{enumerate}

\subsubsection{\'Etape de Modification des Services Existants}

La gestion des modifications de services en cours de vie (r\'esiliation, d\'em\'enagement, upgrade\,/\,downgrade) \'etait particuli\`erement probl\'ematique~:

\begin{enumerate}
    \item L'agent devait identifier manuellement les lignes concern\'ees parmi les abonnements existants~;
    \item Les types de modification (r\'esiliation, remplacement, cr\'eation administrative) n'\'etaient pas encadr\'es par des workflows distincts~;
    \item Les erreurs de saisie sur une ligne \'etaient difficiles \`a identifier, aucun m\'ecanisme de notification granulaire n'existant~;
    \item La tra\c{c}abilit\'e des changements d'\'etat des abonnements \'etait assur\'ee manuellement, avec un risque \'elev\'e de perte d'information.
\end{enumerate}

\subsection{Processus de Gestion des Abonnements}

\subsubsection{Suivi du Cycle de Vie des Abonnements}

Le suivi des abonnements post-vente pr\'esentait \'egalement de nombreuses lacunes~:

\begin{enumerate}
    \item L'\'etat des abonnements (actif, suspendu, r\'esili\'e) n'\'etait pas visible en temps r\'eel depuis l'interface de devis~;
    \item Les modifications d'abonnements d\'eclenchaient des interventions manuelles dans plusieurs syst\`emes~;
    \item L'historique des changements n'\'etait pas consolid\'e, rendant les audits et les contr\^oles difficiles~;
    \item La facturation associ\'ee aux modifications n'\'etait pas automatis\'ee, engendrant des risques d'\'ecarts.
\end{enumerate}

% ========================================
\section{Identification des Points de Douleur}
% ========================================

\subsection{Probl\`emes Li\'es \`a la Cr\'eation et l'\'Edition des Devis}

\subsubsection{Complexit\'e du Catalogue et Lenteur de Recherche}

\begin{itemize}
    \item \textbf{Catalogue trop vaste}~: avec pr\`es de 45\,000 r\'ef\'erences disponibles, la recherche manuelle d'un article adapt\'e repr\'esentait une charge cognitive et temporelle consid\'erable pour les agents commerciaux.
    \item \textbf{Absence de configurateur}~: aucun assistant de type wizard ne guidait le commercial dans la s\'election des offres, l'obligeant \`a construire chaque devis article par article.
    \item \textbf{Gestion manuelle de la hi\'erarchie}~: la structure lignes\,/\,sous-lignes des devis t\'el\'ecom n'\'etait pas support\'ee nativement, n\'ecessitant des manipulations fastidieuses.
\end{itemize}

\subsubsection{Fiabilit\'e des Calculs et des Remises}

\begin{itemize}
    \item \textbf{Calculs de marges non contr\^ol\'es}~: en l'absence de r\`egles m\'etier automatis\'ees, les commerciaux appliquaient des remises pouvant conduire \`a des devis sous-\'evalu\'es, avec un impact direct sur la rentabilit\'e.
    \item \textbf{Absence de validation en temps r\'eel}~: aucun m\'ecanisme d'alerte n'existait pour signaler qu'une remise appliqu\'ee d\'epassait le seuil de marge acceptable.
    \item \textbf{Inconsistance des donn\'ees}~: les prix et quantit\'es saisis manuellement n'\'etaient pas v\'erifi\'es automatiquement, g\'en\'erant r\'eguli\`erement des \'ecarts entre le devis et le bon de commande.
\end{itemize}

\subsection{Probl\`emes Li\'es \`a la Gestion des Configurations Utilisateur}

\subsubsection{Absence de Persistance des Pr\'ef\'erences}

\begin{itemize}
    \item \textbf{Reconfiguration syst\'ematique}~: les agents devaient reconfigurer manuellement l'affichage des colonnes (visibles, masqu\'ees, ordre) \`a chaque nouvelle session, sans m\'emorisation des pr\'ef\'erences individuelles.
    \item \textbf{Productivit\'e r\'eduite}~: cette contrainte r\'ep\'etitive entra\^inait une perte de temps quotidienne non n\'egligeable pour des \'equipes g\'erant de nombreux devis.
    \item \textbf{Exp\'erience utilisateur d\'egrad\'ee}~: l'interface ne s'adaptait pas au profil de l'utilisateur, rendant l'outil moins intuitif qu'attendu.
\end{itemize}

\subsubsection{Performances Insuffisantes sur les Devis Volumineux}

\begin{itemize}
    \item \textbf{Ralentissements importants}~: les devis comportant plusieurs centaines ou milliers de lignes g\'en\'eraient des comportements anormaux de l'interface (chargement infini, blocages de navigation).
    \item \textbf{Absence de virtualisation}~: le rendu de l'int\'egralit\'e des lignes sans m\'ecanisme de virtualisation \'epuisait les ressources du navigateur et d\'egrad\'ait significativement la r\'eactivit\'e de l'application.
\end{itemize}

\subsection{Probl\`emes Li\'es \`a la Gestion des Modifications de Services}

\subsubsection{Opacit\'e des Erreurs et Difficult\'e de D\'ebogage}

\begin{itemize}
    \item \textbf{Messages d'erreur globaux}~: en cas d'\'echec d'une modification de service, l'interface affichait un message g\'en\'erique sans indiquer quelle ligne d'abonnement \'etait concern\'ee, obligeant l'agent \`a inspecter manuellement l'ensemble du devis.
    \item \textbf{Perte des modifications valides}~: une erreur sur une seule ligne annulait l'ensemble des modifications en cours, y compris celles correctement saisies sur les autres lignes.
    \item \textbf{Absence de feedback actionnable}~: les agents ne disposaient pas d'information suffisante pour corriger l'erreur sans l'intervention d'un support technique.
\end{itemize}

\subsubsection{Absence de S\'election Globale sur Filtres}

\begin{itemize}
    \item \textbf{S\'election page par page}~: lors d'op\'erations en masse sur des listes filtr\'ees, les agents ne pouvaient s\'electionner que les lignes visibles sur la page courante, rendant les traitements multi-pages laborieux.
    \item \textbf{Manque d'information contextuelle}~: aucun indicateur ne signalait combien de lignes correspondaient au filtre actif en dehors de la page visible.
\end{itemize}

\subsection{Probl\`emes Transversaux}

\subsubsection{S\'ecurit\'e et Exposition des Interfaces Techniques}

\begin{itemize}
    \item \textbf{Interface Swagger accessible en production}~: la documentation interactive de l'API .NET \'etait accessible dans les environnements de pr\'e-production et de production, exposant la totalit\'e des endpoints et leurs param\`etres \`a tout utilisateur disposant d'un acc\`es r\'eseau, en contradiction avec les exigences de s\'ecurit\'e de l'entreprise.
    \item \textbf{Absence de g\'en\'eration automatis\'ee de la sp\'ecification}~: la sp\'ecification OpenAPI n'\'etait pas g\'en\'er\'ee automatiquement en pipeline, ce qui imposait des mises \`a jour manuelles et introduisait des risques de d\'esynchronisation entre la documentation et l'impl\'ementation r\'eelle.
\end{itemize}

\subsubsection{Qualit\'e du Code et Maintenabilit\'e}

\begin{itemize}
    \item \textbf{Dette technique accumul\'ee}~: l'analyse SonarQube r\'ev\'elait un nombre significatif de code smells, de duplications et de complexit\'es cyclomatiques \'elev\'ees dans les composants Vue.js, rendant la base de code difficile \`a faire \'evoluer sans risque de r\'egression.
    \item \textbf{Pipelines CI instables}~: les pipelines d'int\'egration continue \'echouaient r\'eguli\`erement en raison de probl\`emes de linting, de formatage non conform et de d\'ependances incompatibles, fragilisant la cha\^ine de livraison.
    \item \textbf{D\'ependances obsol\`etes}~: plusieurs biblioth\`eques Python pr\'esentaient des vuln\'erabilit\'es connues ou des versions d\'epass\'ees, sans proc\'edure structur\'ee de mise \`a jour.
\end{itemize}

% ========================================
\section{Impact des Dysfonctionnements}
% ========================================

\subsection{Impact Op\'erationnel}

\subsubsection{Inefficacit\'e des Processus Commerciaux}

Les dysfonctionnements identifi\'es ont g\'en\'er\'e plusieurs impacts n\'egatifs sur l'efficacit\'e des \'equipes~:

\begin{itemize}
    \item \textbf{Perte de temps}~: la recherche manuelle dans un catalogue de 45\,000 r\'ef\'erences et la reconfiguration syst\'ematique de l'interface repr\'esentaient une charge quotidienne injustifi\'ee pour les agents.
    \item \textbf{Cycles de validation prolong\'es}~: l'absence de workflow structur\'e allongeait consid\'erablement les d\'elais de validation et de transmission des devis aux clients.
    \item \textbf{Difficult\'e de traitement en masse}~: l'impossibilit\'e de s\'electionner globalement des lignes filtr\'ees contraignait les agents \`a des op\'erations p\'enibles et r\'ep\'etitives sur les grands portefeuilles de devis.
\end{itemize}

\subsubsection{Qualit\'e de Service D\'egrad\'ee}

\begin{itemize}
    \item \textbf{Exp\'erience utilisateur m\'ediocre}~: les ralentissements sur les devis volumineux, les messages d'erreur opaques et les reconfigurations r\'ep\'etitives nuisaient \`a l'adoption de l'outil par les agents.
    \item \textbf{Frustration des \'equipes}~: la complexit\'e des proc\'edures et l'absence d'automatisations \'el\'ementaires g\'en\'eraient un sentiment de charge inutile chez les utilisateurs.
    \item \textbf{Risque d'erreur humaine \'elev\'e}~: la multiplication des saisies manuelles augmentait m\'ecaniquement la probabilit\'e d'erreurs dans les devis transmis aux clients.
\end{itemize}

\subsection{Impact S\'ecuritaire et Technique}

\subsubsection{Risques de S\'ecurit\'e}

\begin{itemize}
    \item \textbf{Exposition de la documentation API}~: l'accessibilit\'e de Swagger\,UI en environnements sensibles constituait un vecteur de reconnaissance pour tout acteur malveillant disposant d'un acc\`es r\'eseau, en permettant d'inventorier l'ensemble des endpoints, des types de param\`etres et des r\'eponses attendues.
    \item \textbf{Vuln\'erabilit\'es de d\'ependances}~: des biblioth\`eques Python pr\'esentant des CVE connues \'etaient utilis\'ees en production sans processus de mise \`a jour formalis\'e.
\end{itemize}

\subsubsection{Risques de Maintenabilit\'e}

\begin{itemize}
    \item \textbf{Dette technique croissante}~: sans processus de revue de code syst\'ematique ni seuil de qualit\'e appliqu\'e en CI, la base de code se d\'egrad\'ait progressivement, augmentant le co\^ut des \'evolutions futures.
    \item \textbf{Observabilit\'e inexistante}~: l'absence de centralisation des logs applicatifs rendait le diagnostic des incidents en production lent et incomplet, sans possibilit\'e de corr\'elation entre les \'ev\'enements survenus sur les diff\'erentes couches de l'application.
\end{itemize}

\subsection{Impact Strat\'egique}

\subsubsection{Comp\'etitivit\'e et Image}

\begin{itemize}
    \item \textbf{Lenteur des r\'eponses commerciales}~: des cycles de devis prolong\'es r\'eduisaient la r\'eactivit\'e commerciale de l'entreprise face \`a ses concurrents.
    \item \textbf{Image de l'outil}~: une interface peu performante et peu ergonomique nuisait \`a l'image interne de la solution et freinerait son adoption \`a plus grande \'echelle.
    \item \textbf{Difficult\'e d'\'evolution}~: une base de code de faible qualit\'e et des pipelines instables rendaient co\^uteuse l'introduction de nouvelles fonctionnalit\'es, limitant la capacit\'e \`a r\'epondre aux \'evolutions du march\'e.
\end{itemize}

% ========================================
\section{D\'efinition de la Probl\'ematique}
% ========================================

\subsection{Probl\'ematique Principale}

Au regard de l'analyse des processus existants et des dysfonctionnements identifi\'es, la probl\'ematique principale peut \^etre formul\'ee ainsi~:

\begin{quote}
\textit{Comment concevoir et d\'evelopper une plateforme web sp\'ecialis\'ee qui permette de moderniser l'ensemble du processus de gestion des devis et des abonnements dans le secteur des t\'el\'ecommunications B2B~--- en int\'egrant un configurateur d'offres, un moteur de calcul des marges fiable, une gestion granulaire des modifications de services et une exp\'erience utilisateur adapt\'ee aux contraintes des agents commerciaux~--- tout en garantissant la s\'ecurit\'e, la performance et la maintenabilit\'e de la solution dans un environnement distribu\'e et conteneuris\'e~?}
\end{quote}

\subsection{Sous-probl\'ematiques Techniques}

Cette probl\'ematique principale se d\'ecline en plusieurs sous-probl\'ematiques sp\'ecifiques~:

\subsubsection{Conception de l'Architecture et des Flux}

\begin{itemize}
    \item Comment architecturer une solution trois couches (frontend, API, backend m\'etier) capable de s'int\'egrer nativement avec les syst\`emes CRM et ERP existants~?
    \item Comment mod\'eliser les flux d'appels entre les composants (cr\'eation, \'edition, validation, archivage) de mani\`ere \`a garantir la coh\'erence des donn\'ees \`a chaque \'etape~?
    \item Comment g\'erer la s\'eparation entre services \`a vendre, services \`a produire et services \`a exploiter au sein d'une m\^eme structure de devis~?
\end{itemize}

\subsubsection{Performance et Ergonomie}

\begin{itemize}
    \item Comment assurer la fluidit\'e de l'interface pour des devis comportant plusieurs milliers de lignes, tout en maintenant la r\'eactivit\'e des interactions utilisateur~?
    \item Comment concevoir un syst\`eme de persistance des pr\'ef\'erences d'affichage par utilisateur, compatible avec un environnement multi-sessions et multi-tenants~?
    \item Comment r\'eduire la charge cognitive des agents commerciaux en proposant un acc\`es structur\'e au catalogue d'offres, via des configurateurs adapt\'es aux types de services propos\'es~?
\end{itemize}

\subsubsection{S\'ecurit\'e et Observabilit\'e}

\begin{itemize}
    \item Comment contr\^oler finement l'exposition des interfaces techniques (documentation API, endpoints sensibles) selon l'environnement de d\'eploiement~?
    \item Comment mettre en place une strat\'egie de qualit\'e logicielle automatised (linting, formatage, couverture de tests) int\'egr\'ee aux pipelines de livraison continue~?
    \item Comment poser les fondations d'une architecture d'observabilit\'e centr\'ee sur la collecte et l'analyse des logs applicatifs et de s\'ecurit\'e \`a l'\'echelle des environnements Kubernetes~?
\end{itemize}

% ========================================
\section{Contexte du Projet Bifr\"ost}
% ========================================

\subsection{Positionnement Strat\'egique}

\subsubsection{Alignement avec la Strat\'egie de l'Entreprise Cliente}

Le projet Bifr\"ost s'inscrit dans la strat\'egie de transformation digitale de l'entreprise cliente, qui vise \`a~:

\begin{itemize}
    \item Moderniser les outils commerciaux internes pour am\'eliorer la r\'eactivit\'e et la fiabilit\'e des processus de vente~;
    \item R\'eduire la d\'ependance aux manipulations manuelles sources d'erreurs et de lenteurs~;
    \item Am\'eliorer l'exp\'erience des agents du service client, qui constituent les utilisateurs quotidiens de la plateforme~;
    \item Cr\'eer les conditions d'une \'evolution continue de la solution, adapt\'ee aux \'evolutions du march\'e t\'el\'ecom.
\end{itemize}

\subsubsection{R\^ole de VISEO dans le Projet}

VISEO, en tant que soci\'et\'e de conseil et d'int\'egration num\'erique, a jou\'e un r\^ole central dans ce projet en apportant~:

\begin{itemize}
    \item Son expertise en d\'eveloppement full-stack (Vue.js, .NET, Django) et en architecture de syst\`emes distribu\'es~;
    \item Sa ma\^itrise des m\'ethodologies Agile\,/\,Scrum pour structurer la livraison en sprints \'evolutifs~;
    \item Sa capacit\'e \`a concevoir des solutions maintenables, s\'ecuris\'ees et conformes aux standards de qualit\'e logicielle~;
    \item Son exp\'erience dans l'accompagnement de projets complexes n\'ecessitant une int\'egration profonde avec des syst\`emes existants.
\end{itemize}

\subsection{Enjeux et D\'efis du Projet}

\subsubsection{Enjeux Techniques}

\begin{itemize}
    \item \textbf{Int\'egration}~: concevoir une solution s'int\'egrant nativement avec les syst\`emes CRM (Ga\"ia\,/\,Stargate) et ERP existants via des API REST bien d\'efinies.
    \item \textbf{Performance}~: garantir la fluidit\'e de l'interface pour des catalogues volumineux et des devis de grande taille, gr\^ace \`a des m\'ecanismes de virtualisation c\^ot\'e frontend.
    \item \textbf{S\'ecurit\'e}~: contr\^oler finement l'exposition des interfaces selon les environnements et impl\'ementer une authentification robuste via Keycloak.
    \item \textbf{\'Evolutivit\'e}~: adopter une architecture modulaire permettant d'int\'egrer de nouvelles fonctionnalit\'es (configurateur t\'el\'ecom, observabilit\'e) sans refonte majeure.
\end{itemize}

\subsubsection{Enjeux Organisationnels}

\begin{itemize}
    \item \textbf{Adoption utilisateur}~: concevoir une interface suffisamment intuitive pour des agents aux profils vari\'es, minimisant la courbe d'apprentissage~;
    \item \textbf{Processus}~: formaliser et automatiser les workflows m\'etier existants (validation de devis, gestion des abonnements) au sein de la plateforme~;
    \item \textbf{Qualit\'e}~: mettre en place des garde-fous techniques (CI/CD, SonarQube, revues de code) garantissant la p\'erennit\'e de la solution dans le temps.
\end{itemize}

\subsubsection{Enjeux Strat\'egiques}

\begin{itemize}
    \item \textbf{Comp\'etitivit\'e commerciale}~: r\'eduire les cycles de production des devis pour am\'eliorer la r\'eactivit\'e commerciale de l'entreprise cliente face \`a ses concurrents~;
    \item \textbf{R\'eplicabilit\'e}~: cr\'eer une solution architectur\'ee de mani\`ere \`a pouvoir s'adapter \`a d'autres contextes ou clients similaires~;
    \item \textbf{Observabilit\'e}~: poser les bases d'une supervision applicative compl\`ete, pr\'ealable indispensable \`a la mise en conformit\'e avec les exigences de s\'ecurit\'e de l'entreprise.
\end{itemize}

% ========================================
\section*{Conclusion}
% ========================================

L'analyse approfondie des processus existants de gestion des devis et des abonnements t\'el\'ecom au sein de l'entreprise cliente a r\'ev\'el\'e de nombreux dysfonctionnements ayant un impact direct sur l'efficacit\'e op\'erationnelle, la s\'ecurit\'e de l'application et la qualit\'e de l'exp\'erience utilisateur. Ces limitations, h\'erit\'ees d'un outil g\'en\'eraliste inadapt\'e aux sp\'ecificit\'es m\'etier du secteur, ont justifi\'e pleinement la conception d'une solution sur mesure et moderne.

La probl\'ematique identifi\'ee d\'epasse le simple cadre technique pour englober des enjeux organisationnels, s\'ecuritaires et strat\'egiques majeurs. Le projet Bifr\"ost a ainsi repr\'esent\'e une opportunit\'e de transformer en profondeur les pratiques commerciales de l'entreprise cliente, tout en contribuant \`a la strat\'egie globale de modernisation de son syst\`eme d'information.

Les d\'efis techniques et fonctionnels identifi\'es dans ce chapitre ont orient\'e l'ensemble des choix de conception et de d\'eveloppement pr\'esent\'es dans les chapitres suivants, en garantissant que la solution r\'eponde pr\'ecis\'ement aux besoins exprim\'es et aux objectifs strat\'egiques de l'entreprise.